\documentclass[10pt,a4paper]{amsart}
\usepackage[margin=19mm]{geometry}
\usepackage{amsmath,amssymb,mathtools}
\usepackage[scr=rsfs,cal=euler]{mathalfa}
\usepackage{xfrac}
\usepackage[unicode,hidelinks,bookmarks=false]{hyperref}
\newcommand{\ie}{i.e.\ }
\newcommand{\cf}{cf.\ }
\newcommand{\resp}{respectively\ }
\newcommand{\Subsection}{\subsection}
\providecommand{\nobreakdash}{\nobreak\hbox{-}}
\setlength{\emergencystretch}{2em}
\multlinegap0pt
\usepackage{amssymb}
\usepackage{tikz}
\usepackage{mathtools}
\usepackage[shortlabels]{enumitem}
\newtheorem{proposition}{Proposition}
\def\N{{\mathbb N}}
\title{Maximal inequalities, frames and greedy algorithms}
\author{Pablo Bern\'a, Daniel Freeman, Timur Oikhberg, Mitchell Taylor}
\date{Source: arXiv:2601.01047v1 (CC BY 4.0)}
\begin{document}
\maketitle

\noindent\emph{Source and licence.} Maximal inequalities, frames and greedy algorithms. Version arXiv:2601.01047v1 (CC BY 4.0), distributed by arXiv under the Creative Commons Attribution 4.0 International licence, http://creativecommons.org/licenses/by/4.0/, as recorded in the arXiv licence metadata for this exact version and shown on the abstract page https://arxiv.org/abs/2601.01047v1. Full licence notice: \texttt{licenses/arxiv-2601.01047-CC-BY-4.0.txt}.\par\smallskip

\noindent\textit{Editorial scope.} Counted unit: the article's proposition p:destroy\_Lind (source0.tex lines 965--967). The article's complete original proof of it is delivered with it (source0.tex lines 968--1012, one \texttt{proof} environment of 45 lines). No mathematics is written, repaired or re-derived: the statement and the proof are the article's own lines, in the article's own notation, and no retained line is substituted or re-flowed.  No further source-local context is needed: every cross-reference of every retained line resolves inside the two delivered spans.  Nothing was omitted from any retained span, so the package carries no omission record.  No retained line contains a citation, so the wrapper carries no bibliography and no bibliography entry.  No retained line contains a figure environment, an image-inclusion command or drawing code, so no image is delivered and the package declares no asset dependency.  Editorial formatting only, in addition: an AMS \texttt{amsart} preamble that restates the article's own declarations and macros used by the retained lines, namely the environments \texttt{proposition} on separate counters, together with the standard packages those lines need.  The retained environments print in this document as Proposition 1 in order of appearance, whereas the article prints the same items with the numbering of its own full document.  The complete native source of the article as distributed in the arXiv e-print for this exact version is bundled unchanged as \texttt{sources/192045/source0.tex}.
\begin{proposition}\label{p:destroy_Lind}
 In the above notation, the sequence $\big(x_i^{(n)}\big)_{1 \leq i \leq n, n \in \N}$ is neither uniformly quasi-greedy nor bibasic.
\end{proposition}

\begin{proof}
For any $C > 1$, find $m \in \N$ with $m+1 > C/2$.
 For $n \geq 3 \cdot 2^m$, we construct $x \in {\mathrm{span}} \big[ \big(x_i^{(n)}\big)_{1 \leq i \leq n} \big]$ which witnesses the fact that the uniformly quasi-greedy and  bibasis constants of the sequence $\big(x_i^{(n)}\big)_{1 \leq i \leq n}$ in $\ell_1^{2n+2}$ are at least $C$.
 Since $n$ is fixed, we shall use $x_i$ instead of $x_i^{(n)}$.
 \medskip

 
 For $\overline{j} = ( j_1, \ldots, j_k ) \in \{1,2\}^k$ define $\phi_{\overline{j}}(t) = \phi_{j_k} \big( \ldots \phi_{j_2}(\phi_{j_1}(t)) \ldots \big)$; then
 $\phi_{\overline{j}}(t) = 2^k t + 2^{k-1} j_1 + \ldots + 2 j_{k-1} + j_k$.
 For convenience, let $\phi_\emptyset(t) = t$.
 Note that, if $\overline{j} = ( j_1, \ldots, j_k )$ and $\overline{i} = ( i_1, \ldots, i_\ell )$, then $\phi_{\overline{j}}(1) = \phi_{\overline{i}}(1)$ holds iff $\ell = k$ and $j_1=i_1,\dots, j_k=i_k$. Indeed, for $k = \ell$ this follows from the uniqueness of the binary decomposition. On the other hand, if $\ell > k$ then
 $$
 \phi_{\overline{j}}(1) \leq 2^k + 2 \big( 2^{k-1} + \ldots + 1 \big) < 2^k + \big( 2^k + 2^{k-1} + \ldots + 1 \big) \leq \phi_{\overline{i}}(1) .
 $$
 The case of $\ell < k$ is handled in a similar way.
 \medskip
 
 For $m \in \N$, let
 $$
 y_m = \sum_{k=0}^{m-1} 2^{-k} \sum_{j_1, \ldots, j_k \in \{1,2\}} x_{\phi_{j_k, \ldots, j_1}(1)},
 $$
 where, by convention, the term for $k=0$ corresponds to $x_{\phi_\emptyset(1)} = x_1$.
 % Note that the natural ordering of the summands of $x$ coincides with the greedy ordering. Therefore, it suffices to show that $y_m$ (for $m$ large) witnesses the failure of the bibasis property of $(x_i)$.  \medskip
 A direct computation shows that
 \begin{equation}
     \label{eq:y-m1}
 y_m = e_1 - 2^{-m} \sum_{j_1, \ldots, j_m \in \{1,2\}} e_{\phi(j_m, \ldots, j_1)(1)} .
 \end{equation}
 By the definition of $\phi$,
 \begin{equation}
     \label{eq:y-m2}
 \sum_{j_1, \ldots, j_m \in \{1,2\}} e_{\phi(j_m, \ldots, j_1)(1)} = % \sum_{i=2^{m+1}-1}^{2^{m+1}+2^m-2} e_i ,
 \sum_{i \in I_m} e_i, \, {\textrm{  where   }} \, I_m = \{2^{m+1}-1 , \ldots , 2^{m+1}+2^m-2 \} ,
 \end{equation}
 hence $\|y_m\| = 2$.
\medskip

 Next, we show that, for any $N \in \N$, $\| \vee_{m=0}^{N-1} |y_m| \| = N+1$. To this end, observe that, in \eqref{eq:y-m2}, $|I_m| = 2^m$, and $I_m \cap I_k = \emptyset$ if $m \neq k$. Therefore,
 $\vee_{m=0}^{N-1} |y_m| = e_1 + \sum_{m=1}^N 2^{-m} \sum_{i \in I_m} e_i$, yielding the desired estimate.
\medskip

 Note that, if the Lindenstrauss basis is bibasic (uniformly quasi-greedy), then there exists ${\mathbf{B}} \in (0,\infty)$ so that $(x_i)_{i=1}^n$ is bibasic (resp.~uniformly quasi-greedy) with constant ${\mathbf{B}}$, no matter how large $n$ is.
Fix $N$, and pick $n \geq \phi_{2,\ldots,2}(1)$ ($2$ is repeated $N$ times).
If $m \leq N$, then $y_m$ is both a greedy sum and a partial sum of $y_N$, hence $\| \vee_{m=0}^{N-1} |y_m| \| \leq {\mathbf{B}} \|y_N\|$. This gives ${\mathbf{B}} \geq(N+1)/2$ (here ${\mathbf{B}}$ can stand for either the bibasic or the uniformly quasi-greedy constant of $(x_i)_{i=1}^n$), leading to the desired contradiction.
\end{proof}

\end{document}
